% AI-written, AI cross-reviewed; no human sign-off; not a claim that RH is proved
% Atlas unit D174. Local manuscript v3, item G.184, heading on printed PDF p. 203.
% GlobalTheory.T.charpoly_deRhamHomology_endomorphism_eq_map_ladic
% paper2/decomposition_tables.tex:1565–1570. Verbatim excerpt; SHA-256 8b85a0c2826deee789a3b66162dce59dcafbfc0e7cc036e2b21288403898211e.
% Original TeX uses the macros and cross-reference keys of the source manuscript.

\begin{formalstatement}[D174]\label{decomp:D174}
For the given endomorphism, the de Rham homology characteristic polynomial is the image of the integral characteristic polynomial for the \ensuremath{\ell}-adic endomorphism.
\end{formalstatement}
{\small\raggedright\noindent\textsf{\bfseries Related reference}\quad \nolinkurl{GlobalTheory.T.charpoly_deRhamHomology_endomorphism_eq_map_ladic}\par\noindent \nolinkurl{Development/AbelianVarietyTheory/Proofs/DivisorProofs/CharpolyDeRhamHomologyEndomorphismEqMapLadicStatement.lean:8--18}\par\noindent SHA-256 \nolinkurl{97320057fceb6f7ca31418807f0a2f46f59b9c8d15db3ba6bdfca84650a4d128}\par\noindent Dominated 12,706\quad Disjoint 12,706\quad terminal.\par}
The integral characteristic polynomial must be transported to the source’s rational and de Rham homology coordinates.
\hyperref[lit:S040]{Published sources and search S040}.
